\section{Setting and estimators}\label{sec:setup}
\textbf{Estimand.}
A graph on $n$ units has degrees $d_i$ (mean $d$). A design draws the assignment $z\in\{0,1\}^n$, and unit $i$'s outcome $Y_i(z)$ may depend on every coordinate of $z$.
The global average treatment effect is $\gate=\frac1n\sum_i\bigl[Y_i(\mathbf 1)-Y_i(\mathbf 0)\bigr]$, the contrast between everyone treated and everyone in control.
We hold the potential outcomes fixed, noise included, so that $\gate$ is a known number of each simulated population and all randomness of an estimate $\that$ is the assignment.

\textbf{Interference family.}
Let $g_i$ be the fraction of $i$'s neighbours that are treated and $g_i^{(2)}$ the same quantity over two-hop walks. We simulate
\begin{equation}
Y_i(z)=\alpha_i+\beta z_i+\gamma\bigl[(1-w)H(g_i)+w\,H(g_i^{(2)})\bigr]+\sigma\varepsilon_i,\qquad H(g)=(1-\kappa)\,g+\kappa\,h(g),\label{eq:family}
\end{equation}
with $h(0)=0$, $h(1)=1$, so that $\gate=\beta+\gamma$ for every curvature $\kappa\in[0,1]$, range weight $w$ and shape $h$ (convex $g^2$, concave $2g-g^2$, or a steep logistic threshold at $g=\tfrac12$).
We set $\beta=\gamma=1$ ($\gate=2$) unless stated; further variants make the response depend on the treated-neighbour \emph{count} (degree-multiplicative) or let $\beta$ and $\gamma$ vary with degree. The noise $\sigma\varepsilon_i$ is frozen across the two arms; $\sigma=0$ isolates design variance.

\textbf{Designs.}
Bernoulli($p$) assigns units independently. Graph cluster randomisation \citep{ugander2013graph} assigns Louvain communities, with resolution as the cluster scale. A staggered rollout treats an increasing random subset in each of several stages \citep{cortez2022staggered}.

\textbf{Estimators.}
Difference in means (\est{DM}) ignores interference.
A regression estimator maps each unit's assignment to a feature vector $\varphi_i(z)$, fits $y$ on $\varphi$ by least squares with coefficient $\hat b$, and predicts at the two corners:
\begin{equation}
\that=\bar\Delta^{\top}\hat b,\qquad \bar\Delta=\frac1n\sum_i\bigl[\varphi_i(\mathbf 1)-\varphi_i(\mathbf 0)\bigr].\label{eq:reg}
\end{equation}
\est{REG-L} uses $(1,z_i,g_i)$, for which $\that=\hat b_z+\hat b_g$; \est{REG-Q} adds $g_i^2$ and $z_ig_i$; \est{REG-G} adds $g_i^{(2)}$, $z_ig_i$ and log degree; \est{REG-O} uses the true shape of~\eqref{eq:family}, fitted only for its coefficients, and so separates estimation error from misspecification; boosted trees and a graph convolutional network are the learned end of the family.
Design-based comparisons are the Horvitz--Thompson (\est{HT}) and Hájek estimators with the exposure condition ``unit and all neighbours in one arm'' \citep{aronow2017estimating}, and a relaxed condition that accepts a unit when at least $1-\delta$ of its neighbours share its arm.
Staggered rollout (\est{STAG}) observes the all-treated corner but needs several measurements and no carry-over, so it is not like-for-like.

\textbf{Graphs and replicates.}
Synthetic graphs are Erd\H{o}s--R\'enyi (ER), stochastic block model (SBM, ten blocks, 70\% of edges within), Barab\'asi--Albert (BA) and Watts--Strogatz (WS), with $n\approx 2000$ unless stated and mean degree $d$ from 2 to 80.
Two public Facebook networks supply real topology: Stanford3 \citep{traud2012social,rossi2015network} and ego-Facebook \citep{mcauley2012learning}; their edges can be thinned to vary the degree.
A replicate redraws the graph (synthetic only), the noise and the assignment; every cell reports bias, standard deviation and RMSE of $\that-\gate$ with Monte-Carlo intervals, and estimators are compared on the same draws.
